\documentclass{article}
\usepackage{authblk}
\usepackage{amsmath}
\usepackage{biblatex}
\usepackage{graphicx}
\usepackage[hidelinks]{hyperref}
\usepackage[margin=1in]{geometry}
\usepackage{fancyhdr}
\usepackage{rotating}

\makeatletter
\newcommand{\doctitle}[1]{\gdef\@doctitle@text{#1}}
\newcommand{\shorttitle}[1]{\gdef\@shorttitle@text{#1}}
\gdef\@doctitle@text{Untitled}    % defaults so nothing is undefined
\gdef\@shorttitle@text{Untitled}
\newcommand{\thedoctitle}{\@doctitle@text}
\newcommand{\theshorttitle}{\@shorttitle@text}
\makeatother

\setlist[enumerate,1]{label=\arabic*.}
\setlist[enumerate,2]{label=\theenumi\arabic*.}
\setlist[enumerate,3]{label=\theenumii\arabic*.}
\setlist[enumerate,4]{label=\theenumiii\arabic*.}

\pagestyle{fancy}
\fancyhf{} % Clear default headers and footers
\fancyhead[L]{Digital Multimeter Project: \theshorttitle} % Title on the left
\fancyhead[R]{Small et al. \thepage} % Page number on the right
\renewcommand{\headrulewidth}{0pt} % Optional header line

\title{Digital Multimeter:~\thedoctitle}
\date{\today}

\author[1]{Jacob Small}
\author[2]{Jacob Goin}
\author[3]{Micheal Keneipp}

%majors and minors
\affil[1]{Computer Engineering Technology Major, Civic Leadership Minor}
\affil[2]{Computer Engineering Technology Major, CAD Minor}
\affil[3]{Electrical Engineering Technology Major}

\addbibresource{bibliography.bib}
\doctitle{ECT 437 and ECT 438 Project Research and Implementation}
\shorttitle{Research and Implementation}

\begin{document}
    \maketitle
    \section{Overview}
    As discussed in the Project Scope, this project aims to design a digital multimeter (DMM)
    system that can measure current, voltage, resistance, compliance, and inductance.
    To best complete this goal, we need to research the methods used by devices that can measure current, voltage,
    resistance, compacitance, and inductance.

    This document will achieve four main items.
    First, in~\autoref{sec:lit_rev}, we will discuss literature surrounding the concepts and processes needed for
    this project.
    Second, in~\autoref{sec:hard_rev}, we will deconstruct consumer-grade DMMs, and examine how they do the
    measurements that they do, as well as finding items that can be implemented within our system.
    Third, in~\autoref{sec:soft_rev}, we will examine any companion software to the hardware in~
    \autoref{sec:hard_rev} in order to find items that can be implemented within our system.
    Fourth, in~\autoref{sec:conc}, we will discuss how to achieve the goals of the project via the integration of
    both hardware and software.

    \section{Literature Review}\label{sec:lit_rev}
    \autoref{tab:ac_com_dmms} shows the voltage capabilities of common DMMs, including their full scale, accuracy,
    resolution, and frequency range.~\parencite[24]{lecklider2012handheld}
    \begin{sidewaystable}
    \centering
    \caption{Voltage Capabilities of Common DMMs}
    \label{tab:ac_com_dmms}
    \begin{tabular}{cc|ccc|cccc}\hline
        & & \multicolumn{3}{c|}{DC}& \multicolumn{4}{c}{AC}\\\cline{3-9}
        Make& Model& Full Scale& Accuracy& Resolution& Full Scale& Accuracy& Resolution& Frequency\\
        & & $V$& $\pm\%+\text{Counts}$& $\mu V$& $(V\text{ to }V)$& $\pm\%+\text{Counts}$& $\mu V$& $Hz$\\\hline
        AMEC& 5233& $0.06\to1k$& $1\%+12$& $10$& $0.06\to1k$& $1\%+12$& $10$& $10\to3k$\\
        AMEC& MTX3283-BCOMCM& $0.1\to1k$& $0.08\%+8$& $1$& $0.1\to1k$& $0.3\%+40$& $1.0$& $45\to200k$\\\hline
        Agilent Technologies& U1273AX& $0.03\to1k$& $0.05\%+2$& $1$& $0.03\to1k$& $0.6\to20$& $1$& $45\to65$\\
        Agilent Technologies& U1233A& $0.6\to600$& $0.5\%+2$& $100$& $0.6\to600$& $1\%+3$& $100$& $45\to500$\\
        Agilent Technologies& U1253B& $0.05\to1k$& $0.025\%+5$& $1$& $0.05\to1K$& $0.4\%+25$& $1$& $45\to1k$\\\hline
        Amprobe& HD160C& $1\to1.5k$& $0.1\%+5$& $100$& $1\to1K$& $1.2\%+10$& $100$& $45\to2k$\\
        Amprobe& 34SXR-A& $1\to1k$& $0.25\%+5$& $100$& $1\to750$& $1.2\%+10$& $100$& $45\to2k$\\
        Amprobe& 38SXR-A& $0.4\to1k$& $0.5\%+1$& $100$& $0.4\to750$& $1.2\%+8$& $100$& $45\to2k$\\\hline
        B \& K Precision& 2712& $0.4\to1k$& $0.1\%+5$& $10$& $45\to1K$& $1.2\%+20$& $10$& $45\to1k$\\
        B \& K Precision& 2709B& $0.66\to1k$& $0.5\%+2$& $100$& $0.66\to750$& $1.5\%+8$& $100$& $50\to500$\\
        B \& K Precision& 2708B& $0.34\to1k$& $0.1\%+2$& $100$& $3.4\to750$& $1.5\%+8$& $1k$& $50\to1k$\\\hline
        Dranetz-BMI& DronTectl PMIT& $3\to300$& $0.1\%+5$& $10$& $3\to600$& $0.2\%+10$& $100$& $20\to1k$\\
        Dranetz-BMI& DronTectl ISO& $0.3\to1k$& $0.15\%+2$& $100$& $0.3\to1K$& $1\%+3$& $100$& $15\to10k$\\
        Dranetz-BMI& DronTectl Outdoor& $0.1\to1k$& $0.5\%+3$& $10$& $0.1\to1K$& $0.5\%+3$& $10$& $15\to10k$\\\hline
        Extech& EX540& $0.4\to1k$& $0.06\%+2$& $10$& $0.4\to1K$& $0.2\%+10$& $100$& $20\to1k$\\
        Extech& MM57QA& $0.5\to1k$& $0.02\%+2$& $1$& $0.5\to1K$& $1\%+3$& $100$& $15\to10k$\\
        Extech& MG300& $0.4\to1k$& $0.06\%+4$& $10$& $0.4\to1K$& $0.5\%+3$& $10$& $15\to10k$\\\hline
        Fluke& CNX 3000& $0.6\to1k$& $0.09\%+2$& $100$& $0.6\to1K$& $1\%+3$& $100$& $45\to1k$\\
        Fluke& 289& $0.05\to1k$& $0.02\%+2$& $1$& $0.05\to1K$& $0.35\%+25$& $1$& $20\to100k$\\
        Fluke& 28 II& $0.6\to1k$& $0.09\%+1$& $100$& $0.6\to1K$& $0.7\%+2$& $100$& $45\to20k$\\\hline
        Yokogawa USA& TY520& $0.6\to1k$& $0.09\%+2$& $100$& $0.6\to1K$& $0.5\%+5$& $100$& $40\to1k$\\
        Yokogawa USA& TY720& $0.05\to1k$& $0.02\%+2$& $1$& $0.05\to1K$& $0.4\%+30$& $1$& $10\to100k$\\
        Yokogawa USA& CA450 Process DMM& $0.6\to1k$& $0.09\%+2$& $100$& $0.6\to1K$& $0.5\%+5$& $100$& $45\to500$\\\hline
    \end{tabular}
\end{sidewaystable}

    \begin{table}
    \centering
    \caption{General Information of Common DMMs}
    \label{tab:gen_com_dmms}
    \begin{tabular}{cc|cc}\hline
        Make& Model& Cost& Number Sold\\\hline
        AMEC& 5233& $\$99$& $6k$\\
        AMEC& MTX3283-BCOMCM& $\$899$& $100k$\\\hline
        Agilent Technologies& U1273AX& $\$450$& $30k$\\
        Agilent Technologies& U1233A& $\$168$& $6k$\\
        Agilent Technologies& U1253B& $\$459$& $50k$\\\hline
        Amprobe& HD160C& $\$287.96$& $10k$\\
        Amprobe& 34SXR-A& $\$148.76$& $10k$\\
        Amprobe& 38SXR-A& $\$98.36$& $4k$\\\hline
        B \& K Precision& 2712& $\$115$& $40k$\\
        B \& K Precision& 2709B& $\$105$& $6.6k$\\
        B \& K Precision& 2708B& $\$99$& $3.4k$\\\hline
        Dranetz-BMI& DronTectl PMIT& $\$1.65k$& $30k$\\
        Dranetz-BMI& DronTectl ISO& $\$825$& $31k$\\
        Dranetz-BMI& DronTectl Outdoor& $\$535$& $12k$\\\hline
        Extech& EX540& $\$299.99$& $40k$\\
        Extech& MM570A& $\$329.99$& $500k$\\
        Extech& MG300& $\$429.99$& $40k$\\\hline
        Fluke& CNX 3000& Ask& $6k$\\
        Fluke& 289& $\$588.50$& $50k$\\
        Fluke& 28 II& $\$460.95$& $20k$\\\hline
        Yokogawa USA& TY520& $\$250$& $6k$\\
        Yokogawa USA& TY720& $\$490$& $50k$\\
        Yokogawa USA& CA450 Process DMM& $\$845$& $33k$\\\hline
    \end{tabular}
\end{table}

    Of the DMMs discussed by~\citeauthor{lecklider2012handheld} and reflected in~\autoref{tab:gen_com_dmms}, the
    Extech MM57QA was the most popular, with $500,000$ units sold at a cost\footnote{The costs reported in~
    \autoref{tab:gen_com_dmms} reflect those listed in the source~\parencite[25]{lecklider2012handheld}.
    These costs are from 2012, and a $45.90\%$ increase is needed to calculate 2026 prices.}
    of $\$329.99$, with the average selling $50,373.91$ units at a rate of $\$437.98$.

    The meters discussed span an understandably wide range in both price and features.
    Nearly all of them offer true-rms measurement, continuity and diode testing, and safety ratings of CAT III $1,000$ V or CAT IV $600$ V,
    usually with CE and often UL, CSA, or IEC 61010-1 certification.
    At the budget end, the meters offer little beyond practical conveniences such as an LCD backlight, auto power
    off, a magnetic hanging strap, AC/DC measurement, inductance, dBm, and duty cycle.
    Others in the general price range also add a bar-graph display and non-contact voltage detection, and a flashlight.

    The mid-range meters mostly distinguish themselves through ruggedness, safety design, and connectivity.
    Some are IP67 rated, and withstand $8kV$ spikes.
    Additionally, some are IP67 waterproof with wireless PC connectivity and $9,999$ record storage.
    This price area also have terminal shutters that block incorrect lead connections, and units add dB math and
    duty cycle.

    At the top end are the specialized and data-focused instruments.
    These can include Bluetooth, a graphical display, crest factor, and pulse counting, graphing, a real-time
    clock, an AC filter, and a 1/4 VGA display Others are able to read wirelessly fro  up to three remote modules at
    once.

    Given the popularity of the Extech MM57QA, we will focus on its construction, features, and components via it's
    user manual.

    \subsection{Components}
    The three main visual component groups are~\parencite[4]{extech_mm57qa}:
    \begin{enumerate}
        \item A $5\frac{4}{5}$ digits $500k$ counts LCD display.
        \item Function buttons
        \item $3$ Input Jacks
    \end{enumerate}
    \section{Hardware Review}\label{sec:hard_rev}
    \section{Software Review}\label{sec:soft_rev}
    \section{Conclusion}\label{sec:conc}
    \printbibliography
\end{document}